% from main.tex
%%%%%%%%%%  scan idx 243  |  folio 230  %%%%%%%%%%
pseudo-cohérent équivaut à dire que \(f\) est localement de présentation finie au sens de (EGA IV 1.4).
\end{enumerate}

\textbf{Exemples 1.3.} Un morphisme lisse, une immersion régulière (1.1.2) sont des morphismes pseudo-cohérents. Il en est de même d'un morphisme localement de type fini quelconque \(f:X\to Y\) si \(Y\) est localement noethérien.

\textbf{Proposition 1.4.} Soit \(f:X\to Y\) un morphisme localement de type fini, et soit
\[
E\longrightarrow F\longrightarrow G\longrightarrow E[1]
\]
un triangle distingué de \(D(X)\). On a alors le tableau suivant, où \(n\in\mathbb Z\cup\{-\infty\}\) indique un degré de pseudo-cohérence par rapport à \(f\), et chaque ligne représente une implication dont le but est souligné :
\[
\begin{array}{ccc}
E&F&G\\
n+1&n&\underline n\\
n&\underline n&n\\
\underline n&n&n-1.
\end{array}
\]
En particulier, \(D(f)_{\mathrm{coh}}\) est une sous-catégorie triangulée de \(D(X)\).

\emph{Preuve.} En vertu des propriétés (i) à (iii) de (1.1), on peut supposer que \(f\) est une immersion fermée. On a alors un triangle distingué
\[
f_*(E)\longrightarrow f_*(F)\longrightarrow f_*(G)\longrightarrow f_*(E)[1]
\]
et l'assertion résulte de (I 2.6).

\textbf{Définition 1.5.} Soient \(X\) et \(Y\) deux \(S\)-schémas. Nous dirons que \(X\) et \(Y\) sont tor-indépendants sur \(S\) si l'on a
\[
  \Tor_i^S(\mathcal O_X,\mathcal O_Y)=0\quad\text{pour }i\ne 0,
\]
les \(\Tor_i^S\) étant les hyper-tor locaux définis dans (EGA III 6.5). C'est une
